\documentclass[10pt,a4paper]{amsart}
\usepackage[margin=19mm]{geometry}
\usepackage{amsmath,amssymb,mathtools}
\usepackage[scr=rsfs,cal=euler]{mathalfa}
\usepackage{xfrac}
\usepackage[unicode,hidelinks,bookmarks=false]{hyperref}
\newcommand{\ie}{i.e.\ }
\newcommand{\cf}{cf.\ }
\newcommand{\resp}{respectively\ }
\newcommand{\Subsection}{\subsection}
\providecommand{\nobreakdash}{\nobreak\hbox{-}}
\setlength{\emergencystretch}{2em}
\multlinegap0pt
\usepackage{amssymb}
\usepackage[shortlabels]{enumitem}
\newtheorem{lem}{Lemma}
\newtheorem{prop}{Proposition}
\newcommand{\R}{\mathbb{R}}
\newcommand{\bu}{\boldsymbol{u}}
\newcommand{\iantipodal}{symmetric}
\title{Antipodality of Spherical Designs with Odd Harmonic Indices}
\author{Ryutaro Misawa, Akihiro Munemasa, Masanori Sawa}
\date{Source: arXiv:2410.09471v6 (CC BY 4.0)}
\begin{document}
\maketitle

\noindent\emph{Source and licence.} Antipodality of Spherical Designs with Odd Harmonic Indices. Version arXiv:2410.09471v6 (CC BY 4.0), distributed by arXiv under the Creative Commons Attribution 4.0 International licence, http://creativecommons.org/licenses/by/4.0/, as recorded in the arXiv licence metadata for this exact version and shown on the abstract page https://arxiv.org/abs/2410.09471v6. Full licence notice: \texttt{licenses/arxiv-2410.09471-CC-BY-4.0.txt}.\par\smallskip

\noindent\textit{Editorial scope.} Counted unit: the article's prop prop:antipodal2 (source0.tex lines 186--194). The article's complete original proof of it is delivered with it (source0.tex lines 195--260, one \texttt{proof} environment of 66 lines). No mathematics is written, repaired or re-derived: the statement and the proof are the article's own lines, in the article's own notation, and no retained line is substituted or re-flowed.  Retained uncounted as the source-local context the proof needs: lines 103--117; lines 121--134.  Nothing was omitted from any retained span, so the package carries no omission record.  No retained line contains a citation, so the wrapper carries no bibliography and no bibliography entry.  No retained line contains a figure environment, an image-inclusion command or drawing code, so no image is delivered and the package declares no asset dependency.  Editorial formatting only, in addition: an AMS \texttt{amsart} preamble that restates the article's own declarations and macros used by the retained lines, namely the environments \texttt{lem}, \texttt{prop} on separate counters, together with the standard packages those lines need.  The retained environments print in this document as Lemma 1, Proposition 1 in order of appearance, whereas the article prints the same items with the numbering of its own full document.  The complete native source of the article as distributed in the arXiv e-print for this exact version is bundled unchanged as \texttt{sources/166730/source0.tex}.
\begin{lem}
\label{lem:Newton}
\begin{enumerate}
\item\label{Ni1} 
For all integers \( n \ge k \ge 1 \), we have
\[
k\, e_k(x_1,\ldots,x_n) = \sum_{i=1}^k (-1)^{i-1} e_{k-i}(x_1,\ldots,x_n) \, p_i(x_1,\ldots,x_n).
\]
\item\label{Ni2} 
For all integers \( k \ge n \ge 1 \), we have
\[
0 = \sum_{i=k-n}^k (-1)^{i-1} e_{k-i}(x_1,\ldots,x_n) \, p_i(x_1,\ldots,x_n).
\]
\end{enumerate}
\end{lem}

\begin{prop}
\label{prop:Newton1}
Let \( m \) and \( n \) be positive integers with \( 2m-1 \le n \), and let \( x_1,\ldots,x_n \in \R \). Then the following are equivalent:
\begin{enumerate}
\item\label{N1i1} 
\[
p_{2k-1}(x_1,\ldots,x_n) = 0 \quad \text{for all } k\in\{1,\dots,m\}.
\]
\item\label{N1i2}
\[
e_{2k-1}(x_1,\ldots,x_n) = 0 \quad \text{for all } k\in\{1,\dots,m\}.
\]
\end{enumerate}
\end{prop}

\begin{prop}
\label{prop:antipodal2}
Let \( m \) and \( n \) be positive integers with \( n\le 2m \). Suppose that real numbers \( x_1,\dots,x_n\in[-1,1] \) satisfy
\begin{equation}\label{xi2k-1}
p_{2k-1}(x_1,\dots,x_n)=0
\end{equation}
for all \( k\in\{1,\dots,m\} \). Then the multiset \(\{x_1,\dots,x_n\}\) is \iantipodal. In particular, if a subset \( X\subset[-1,1] \) is an interval \( T_m \)-design with \(|X|=n\), 
then \( X \) is \iantipodal.
\end{prop}

\begin{proof}
Without loss of generality, we may assume \(0\notin X\), because removing zeros does not affect the conditions \(\sum_{x\in X} x^{2k-1}=0\) for \(1\le k\le m\).
We first prove that \eqref{xi2k-1} holds for every positive integer \( k \) by induction on \( k \). Since~\ref{xi2k-1} is assumed for \( k\le m \), so suppose \( k>m \). Then \( 2k-1\ge n \). 
Moreover, because there are only \( n \) variables, \( e_j(x_1,\dots,x_n)=0 \) for all \( j>n \); in particular, for \( j>2m \). 
It follows from Lemma~\ref{lem:Newton}~\ref{Ni2} that
\begin{align*}
p_{2k-1}(x_1,\dots,x_n) &= \sum_{i=2k-1-n}^{2k-2} (-1)^i e_{2k-1-i}(x_1,\dots,x_n) \, p_i(x_1,\dots,x_n)\\
&=\sum_{i=1}^{m} e_{2i-1}(x_1,\dots,x_n) \, p_{2(k-i)}(x_1,\dots,x_n)\\
&\quad-\sum_{i=1}^{m} e_{2i}(x_1,\dots,x_n) \, p_{2(k-i)-1}(x_1,\dots,x_n).
\end{align*}


By Proposition~\ref{prop:Newton1} and \eqref{xi2k-1}, 
we have $e_{2i-1}(x_1,\dots,x_n)=0$ for $1\leq i\leq m$.
Thus, the first term is zero.
By the inductive hypothesis,
the second term
%% the second terms
is also zero.
Therefore, we have shown that
\eqref{xi2k-1} holds for every positive integer $k$.

We now prove the assertion by induction on $n$.
The cases $n=1,2$ are obvious.
Let $Y=\{y_1,\dots,y_{n'}\}$ be the set of all distinct elements of 
%in the special case where 
$x_1,\dots,x_n$, that is, $Y=\{x_1,\dots,x_n\}$ and
$|Y|=n'\leq n$.
%are pairwise distinct. We use induction on $n$.
Define $A\in\R^{n'\times n'}$ by
\[
A_{i,j}=y_j^{2i-1}\quad(1\leq i,j\leq n').
\]
Define a column vector $\bu$ whose $j$-th entry $u_j$ is given by
\[
u_j=|\{i\mid 1\leq i\leq n,\;x_i=y_j\}|\quad(1\leq j\leq n').
\]
Then the $k$-th entry of $A\bu$ is
\[
\sum_{j=1}^{n'} y_j^{2k-1}u_j=p_{2k-1}(x_1,\dots,x_n)=0,
\]
since we have shown that \eqref{xi2k-1} holds for every positive integer $k$.
Thus $A\bu=0$, and hence
%By the assumption, we have $A\bu=0$, where
%$\bu$ is the column vector whose $j$-th entry is 
%and hence
\[
0=\det A=\prod_{k=1}^{n'}y_{k}\prod_{1\leq i<j\leq n'}(y_j^2-y_i^2).
\]
This implies that, either there exists $j$ such that $y_{j}=0$, or
there exist $i,j$ with $i<j$ such that $y_i=-y_j$.
In the latter case, we may assume without loss of generality $x_{n-1}=-x_n$.
%, and
%$X'=X\setminus\{x_i,x_j\}$ in the latter case.
Since we have shown that \eqref{xi2k-1} holds for every positive integer $k$,
we obtain
\[p_{2k-1}(x_1,\dots,x_{n-2})=p_{2k-1}(x_1,\dots,x_{n})=0\]
for every positive integer $k$.
%\sum_{x\in X'}x^{2k-1}=\sum_{x\in X}x^{2k-1}=0
%\quad(1\leq k\leq n),\]
%we have, in particular,
%\[\sum_{x\in X'}x^{2k-1}=\sum_{x\in X}x^{2k-1}=0\quad(1\leq k\leq |X'|).\]
By induction, the multiset $\{x_1,\dots,x_{n-2}\}$ is \iantipodal,
and so is $\{x_1,\dots,x_{n}\}$.
The former case can be proved in a similar manner.
\end{proof}

\end{document}
